\documentclass[12pt,aspectratio=1610]{beamer}
\geometry{paperwidth=20cm,paperheight=12.5cm}
\usepackage[italian,provide=*]{babel}
\usepackage{lmodern,fontspec}
\setsansfont{NimbusSans-Regular.otf}[Path=sorgenti/font/,BoldFont=NimbusSans-Bold.otf,ItalicFont=NimbusSans-Italic.otf]
\setmainfont{NimbusSans-Regular.otf}[Path=sorgenti/font/,BoldFont=NimbusSans-Bold.otf,ItalicFont=NimbusSans-Italic.otf]
\setmonofont{lmmono10-regular.otf}[Path=sorgenti/font/,Scale=.88]
\usefonttheme{professionalfonts}
\usepackage{amsmath,amssymb,booktabs,array,tabularx,graphicx,tikz,listings,bookmark}
\usetikzlibrary{calc,arrows.meta,positioning}
\definecolor{Paper}{HTML}{F8F6F1}
\definecolor{Ink}{HTML}{253344}\definecolor{Navy}{HTML}{15283E}
\definecolor{Quantum}{HTML}{5146A5}\definecolor{Classical}{HTML}{087C74}
\definecolor{Measure}{HTML}{A05B10}\definecolor{Rose}{HTML}{9B4564}
\definecolor{Muted}{HTML}{657180}\definecolor{Rule}{HTML}{D6DAD9}
\definecolor{Wash}{HTML}{EAE7F6}\definecolor{Night}{HTML}{112236}
\definecolor{Lilac}{HTML}{ABB5FF}\definecolor{Mint}{HTML}{72D8C8}\definecolor{Gold}{HTML}{EEC184}
\setbeamercolor{background canvas}{bg=Paper}\setbeamercolor{normal text}{fg=Ink}
\setbeamertemplate{navigation symbols}{}\setbeamertemplate{headline}{}\setbeamertemplate{footline}{}
\setbeamersize{text margin left=0pt,text margin right=0pt}\setbeameroption{hide notes}
\newsavebox{\BodyBox}\newsavebox{\SideBox}\newsavebox{\RecallBox}\newsavebox{\TitleBox}
\newcommand{\MeasureBox}[2]{\typeout{GEOBOX:#1:\the\wd#2:\the\ht#2:\the\dp#2}}
\lstset{language=R,basicstyle=\ttfamily\fontsize{10.2}{13.4}\selectfont,
columns=fullflexible,keepspaces=true,showstringspaces=false,breaklines=true,breakatwhitespace=false,
commentstyle=\color{Muted},keywordstyle=\color{Classical},stringstyle=\color{Quantum},
aboveskip=0pt,belowskip=0pt,xleftmargin=0pt,xrightmargin=0pt,tabsize=2}
\hypersetup{colorlinks=true,linkcolor=Classical,urlcolor=Classical,pdfpagemode=UseOutlines,pdfstartview=Fit,bookmarksopen=true,bookmarksopenlevel=0}
\bookmarksetup{open,openlevel=0,depth=2}
\tikzset{arr/.style={-{Stealth[length=2mm]},line width=.8pt,draw=Classical}}
\hypersetup{pdftitle={Campionamento, stima e incertezza},pdfauthor={},pdfsubject={Corso di statistica e data science per Geologia, primo anno; modulo 05}}
\usetikzlibrary{decorations.pathreplacing,fit,backgrounds,matrix}\hypersetup{pdftitle={P04 · Un geyser: associazione e previsione},pdfsubject={Revisione integrale R6}}\begin{document}
\begin{frame}[plain]\hypertarget{slide1}{}\begin{tikzpicture}[remember picture,overlay]\fill[Night] (current page.south west) rectangle (current page.north east);\draw[Mint,line width=2pt] ([xshift=1.05cm,yshift=-1.35cm]current page.north west)--([xshift=3cm,yshift=-1.35cm]current page.north west);\node[anchor=north west,inner sep=0] at ([xshift=1.05cm,yshift=-1.7cm]current page.north west){\color{Mint}\bfseries\fontsize{11}{14}\selectfont P04 · REVISIONE R6};
\node[anchor=north west,inner sep=0] at ([xshift=1.05cm,yshift=-2.65cm]current page.north west){\begin{minipage}{15.8cm}\raggedright\color{white}\bfseries\fontsize{31}{38}\selectfont Un geyser: associazione e previsione\end{minipage}};
\node[anchor=north west,inner sep=0] at ([xshift=1.05cm,yshift=-6.4cm]current page.north west){\begin{minipage}{15.8cm}\raggedright\color{Lilac}\fontsize{15}{21}\selectfont Durata dell’eruzione, intervalli e valutazione su eventi separati\end{minipage}};
\node[anchor=north west,inner sep=0] at ([xshift=1.05cm,yshift=-8.6cm]current page.north west){\begin{minipage}{16.5cm}\color{white}\fontsize{11}{15}\selectfont 21 slide · teoria, derivazioni, figure e applicazioni in R.\par\smallskip Provenienza dei dati e sorgenti nei materiali allegati.\end{minipage}};
\node[anchor=south east,inner sep=0] at ([xshift=18.95cm,yshift=-11.35cm]current page.north west){\color{Lilac!45}\bfseries\fontsize{64}{68}\selectfont P04};
\end{tikzpicture}\end{frame}
\begin{frame}[plain]\hypertarget{slide2}{}\bookmark[level=0,page=2]{Indice}\begin{tikzpicture}[remember picture,overlay]\node[anchor=north west,inner sep=0] at ([xshift=0.95cm,yshift=-0.65cm]current page.north west){\color{Navy}\bfseries\fontsize{22}{27}\selectfont Percorso del modulo};
\draw[Classical,line width=1.3pt] ([xshift=.95cm,yshift=-1.68cm]current page.north west)--([xshift=2.25cm,yshift=-1.68cm]current page.north west);\node[anchor=north west,inner sep=0] at ([xshift=0.95cm,yshift=-2.23cm]current page.north west){\color{Classical}\bfseries\fontsize{13}{16}\selectfont 3};
\node[anchor=north west,inner sep=0] at ([xshift=2.0cm,yshift=-2.23cm]current page.north west){\begin{minipage}{7.25cm}\raggedright\color{Navy}\bfseries\fontsize{12.4}{15.5}\selectfont\hyperlink{slide3}{Fonte, bersaglio e separazione dei dati}\end{minipage}};
\node[anchor=north west,inner sep=0] at ([xshift=0.95cm,yshift=-3.5300000000000002cm]current page.north west){\color{Classical}\bfseries\fontsize{13}{16}\selectfont 7};
\node[anchor=north west,inner sep=0] at ([xshift=2.0cm,yshift=-3.5300000000000002cm]current page.north west){\begin{minipage}{7.25cm}\raggedright\color{Navy}\bfseries\fontsize{12.4}{15.5}\selectfont\hyperlink{slide7}{Associazione nel campione di adattamento}\end{minipage}};
\node[anchor=north west,inner sep=0] at ([xshift=0.95cm,yshift=-4.83cm]current page.north west){\color{Classical}\bfseries\fontsize{13}{16}\selectfont 12};
\node[anchor=north west,inner sep=0] at ([xshift=2.0cm,yshift=-4.83cm]current page.north west){\begin{minipage}{7.25cm}\raggedright\color{Navy}\bfseries\fontsize{12.4}{15.5}\selectfont\hyperlink{slide12}{Incertezza della media e di un evento}\end{minipage}};
\node[anchor=north west,inner sep=0] at ([xshift=10.299999999999999cm,yshift=-2.23cm]current page.north west){\color{Classical}\bfseries\fontsize{13}{16}\selectfont 15};
\node[anchor=north west,inner sep=0] at ([xshift=11.35cm,yshift=-2.23cm]current page.north west){\begin{minipage}{7.25cm}\raggedright\color{Navy}\bfseries\fontsize{12.4}{15.5}\selectfont\hyperlink{slide15}{Valutazione su eventi separati}\end{minipage}};
\node[anchor=north west,inner sep=0] at ([xshift=10.299999999999999cm,yshift=-3.5300000000000002cm]current page.north west){\color{Classical}\bfseries\fontsize{13}{16}\selectfont 20};
\node[anchor=north west,inner sep=0] at ([xshift=11.35cm,yshift=-3.5300000000000002cm]current page.north west){\begin{minipage}{7.25cm}\raggedright\color{Navy}\bfseries\fontsize{12.4}{15.5}\selectfont\hyperlink{slide20}{Struttura interna e discussione}\end{minipage}};
\node[anchor=north west,inner sep=0] at ([xshift=0.95cm,yshift=-10.55cm]current page.north west){\begin{minipage}{18cm}\color{Muted}\fontsize{10}{13}\selectfont I numeri indicano le pagine iniziali; i titoli sono cliccabili. Codice, figure e verifiche sono inclusi nel progetto.\end{minipage}};
\end{tikzpicture}\end{frame}
\begin{lrbox}{\TitleBox}\begin{minipage}[t]{18.2cm}\raggedright\sloppy\color{Ink}\fontsize{20.7}{24}\selectfont \bfseries\color{Navy}Definire il bersaglio prima del modello\end{minipage}\end{lrbox}
\begin{lrbox}{\BodyBox}\begin{minipage}[t]{10.25cm}\raggedright\sloppy\color{Ink}\fontsize{12.4}{16.4}\selectfont Il dataset faithful contiene 272 coppie di durata dell’eruzione e attesa alla successiva, espresse in minuti. Il progetto esamina se la durata migliori la previsione dell’attesa rispetto a una previsione costante. La durata completa è disponibile soltanto al termine dell’eruzione, mentre l’estratto non fornisce date e protocollo dettagliato dell’origine dei tempi. Un’applicazione operativa dovrebbe quindi verificare anche se l’attesa sia misurata dall’inizio o dalla fine dell’evento, traducendo il bersaglio nel tempo residuo effettivamente disponibile al momento della previsione.\end{minipage}\end{lrbox}
\begin{lrbox}{\SideBox}\begin{minipage}[t]{7.45cm}\raggedright\sloppy\color{Ink}\fontsize{12.2}{16.1}\selectfont \input{tikz/P04/D01.tex}\end{minipage}\end{lrbox}
\begin{lrbox}{\RecallBox}\begin{minipage}[t]{16.45cm}\raggedright\sloppy\color{Ink}\fontsize{10.1}{13.2}\selectfont Il bersaglio operativo richiede un’origine temporale coerente con l’informazione disponibile.\end{minipage}\end{lrbox}
\MeasureBox{3-TitleBox}{\TitleBox}
\MeasureBox{3-BodyBox}{\BodyBox}
\MeasureBox{3-SideBox}{\SideBox}
\MeasureBox{3-RecallBox}{\RecallBox}
\begin{frame}[plain]\hypertarget{slide3}{}
\bookmark[level=0,page=3]{Fonte, bersaglio e separazione dei dati}
\bookmark[level=1,page=3]{Definire il bersaglio prima del modello}
\begin{tikzpicture}[remember picture,overlay]\node[anchor=north east,inner sep=0] at ([xshift=19.05cm,yshift=-0.16cm]current page.north west){\color{Classical}\fontsize{8}{10}\selectfont\hyperlink{slide2}{Indice}};
\node[anchor=north west,inner sep=0] at ([xshift=0.95cm,yshift=-0.62cm]current page.north west){\usebox{\TitleBox}};
\draw[Classical,line width=1.3pt] ([xshift=.95cm,yshift=-1.68cm]current page.north west)--([xshift=2.25cm,yshift=-1.68cm]current page.north west);\node[anchor=north west,inner sep=0] at ([xshift=0.95cm,yshift=-2.12cm]current page.north west){\usebox{\BodyBox}};
\node[anchor=north west,inner sep=0] at ([xshift=11.6cm,yshift=-2.1cm]current page.north west){\color{Quantum}\bfseries\fontsize{11.5}{14}\selectfont Lettura del diagramma};
\node[anchor=north west,inner sep=0] at ([xshift=11.6cm,yshift=-2.75cm]current page.north west){\usebox{\SideBox}};
\draw[Rule,line width=.5pt] ([xshift=.95cm,yshift=-10.43cm]current page.north west)--([xshift=19.05cm,yshift=-10.43cm]current page.north west);\node[anchor=north west,inner sep=0] at ([xshift=0.95cm,yshift=-10.66cm]current page.north west){\color{Classical}\bfseries\fontsize{9.4}{12}\selectfont Richiami};
\node[anchor=north west,inner sep=0] at ([xshift=2.6cm,yshift=-10.64cm]current page.north west){\usebox{\RecallBox}};
\draw[Rule,line width=.35pt] ([xshift=.95cm,yshift=-11.92cm]current page.north west)--([xshift=19.05cm,yshift=-11.92cm]current page.north west);\node[anchor=north west,inner sep=0] at ([xshift=0.95cm,yshift=-12.06cm]current page.north west){\color{Muted}\fontsize{7.7}{9.3}\selectfont P04 · Un geyser: associazione e previsione};
\node[anchor=north east,inner sep=0] at ([xshift=17.65cm,yshift=-12.06cm]current page.north west){\color{Classical}\fontsize{7.3}{9.1}\selectfont R datasets: faithful; Azzalini e Bowman, 1990};
\node[anchor=north east,inner sep=0] at ([xshift=19.05cm,yshift=-12.06cm]current page.north west){\color{Muted}\fontsize{7.7}{9.3}\selectfont 3/21};
\end{tikzpicture}\end{frame}
\begin{lrbox}{\TitleBox}\begin{minipage}[t]{18.2cm}\raggedright\sloppy\color{Ink}\fontsize{20.7}{24}\selectfont \bfseries\color{Navy}La fonte conserva arrotondamenti e limiti storici\end{minipage}\end{lrbox}
\begin{lrbox}{\BodyBox}\begin{minipage}[t]{10.25cm}\raggedright\sloppy\color{Ink}\fontsize{12.4}{16.4}\selectfont La tabella riproduce i valori distribuiti con R e conserva un identificativo di riga, costruito per rendere ripetibile la separazione dei dati. La documentazione segnala arrotondamenti nei tempi, per cui coppie numericamente uguali possono rappresentare eventi distinti. Il controllo confronta i valori con faithful e verifica positività e identificativi, evitando cancellazioni automatiche. L’assenza delle date impedisce di costruire una vera verifica prospettica della stabilità nel tempo e limita l’interpretazione del test a una generalizzazione interna sotto scambiabilità.\end{minipage}\end{lrbox}
\begin{lrbox}{\SideBox}\begin{minipage}[t]{7.45cm}\raggedright\sloppy\color{Ink}\fontsize{12.2}{16.1}\selectfont Il dataset storico riguarda Old Faithful nello Yellowstone. Le unità in minuti e le convenzioni della fonte devono essere mantenute nel modello. Un risultato ottenuto su questo estratto richiederebbe nuove osservazioni cronologiche e verifiche del protocollo prima di sostenere una prestazione operativa attuale.\end{minipage}\end{lrbox}
\begin{lrbox}{\RecallBox}\begin{minipage}[t]{16.45cm}\raggedright\sloppy\color{Ink}\fontsize{10.1}{13.2}\selectfont L’uguaglianza di valori arrotondati va distinta dalla duplicazione accidentale di una registrazione.\end{minipage}\end{lrbox}
\MeasureBox{4-TitleBox}{\TitleBox}
\MeasureBox{4-BodyBox}{\BodyBox}
\MeasureBox{4-SideBox}{\SideBox}
\MeasureBox{4-RecallBox}{\RecallBox}
\begin{frame}[plain]\hypertarget{slide4}{}
\bookmark[level=1,page=4]{La fonte conserva arrotondamenti e limiti storici}
\begin{tikzpicture}[remember picture,overlay]\node[anchor=north east,inner sep=0] at ([xshift=19.05cm,yshift=-0.16cm]current page.north west){\color{Classical}\fontsize{8}{10}\selectfont\hyperlink{slide2}{Indice}};
\node[anchor=north west,inner sep=0] at ([xshift=0.95cm,yshift=-0.62cm]current page.north west){\usebox{\TitleBox}};
\draw[Classical,line width=1.3pt] ([xshift=.95cm,yshift=-1.68cm]current page.north west)--([xshift=2.25cm,yshift=-1.68cm]current page.north west);\node[anchor=north west,inner sep=0] at ([xshift=0.95cm,yshift=-2.12cm]current page.north west){\usebox{\BodyBox}};
\node[anchor=north west,inner sep=0] at ([xshift=11.6cm,yshift=-2.1cm]current page.north west){\color{Quantum}\bfseries\fontsize{11.5}{14}\selectfont Sviluppo e interpretazione};
\node[anchor=north west,inner sep=0] at ([xshift=11.6cm,yshift=-2.75cm]current page.north west){\usebox{\SideBox}};
\draw[Rule,line width=.5pt] ([xshift=.95cm,yshift=-10.43cm]current page.north west)--([xshift=19.05cm,yshift=-10.43cm]current page.north west);\node[anchor=north west,inner sep=0] at ([xshift=0.95cm,yshift=-10.66cm]current page.north west){\color{Classical}\bfseries\fontsize{9.4}{12}\selectfont Richiami};
\node[anchor=north west,inner sep=0] at ([xshift=2.6cm,yshift=-10.64cm]current page.north west){\usebox{\RecallBox}};
\draw[Rule,line width=.35pt] ([xshift=.95cm,yshift=-11.92cm]current page.north west)--([xshift=19.05cm,yshift=-11.92cm]current page.north west);\node[anchor=north west,inner sep=0] at ([xshift=0.95cm,yshift=-12.06cm]current page.north west){\color{Muted}\fontsize{7.7}{9.3}\selectfont P04 · Un geyser: associazione e previsione};
\node[anchor=north east,inner sep=0] at ([xshift=17.65cm,yshift=-12.06cm]current page.north west){\color{Classical}\fontsize{7.3}{9.1}\selectfont };
\node[anchor=north east,inner sep=0] at ([xshift=19.05cm,yshift=-12.06cm]current page.north west){\color{Muted}\fontsize{7.7}{9.3}\selectfont 4/21};
\end{tikzpicture}\end{frame}
\begin{lrbox}{\TitleBox}\begin{minipage}[t]{18.2cm}\raggedright\sloppy\color{Ink}\fontsize{20.7}{24}\selectfont \bfseries\color{Navy}Una separazione condivisa prima dell’esplorazione\end{minipage}\end{lrbox}
\begin{lrbox}{\BodyBox}\begin{minipage}[t]{10.25cm}\raggedright\sloppy\color{Ink}\fontsize{12.4}{16.4}\selectfont Il piano conservato assegna 190 eventi all’insieme di adattamento e 82 a quello di verifica, mantenendo gli identificativi delle righe. L’esplorazione, la stima dei coefficienti e la costruzione della previsione costante utilizzano soltanto i 190 eventi selezionati per l’adattamento. Gli 82 eventi separati vengono impiegati per valutare le procedure già definite. Poiché la divisione è casuale e manca una cronologia, l’esperimento valuta la generalizzazione entro una popolazione di eventi considerata scambiabile, senza simulare necessariamente le condizioni di una previsione futura.\end{minipage}\end{lrbox}
\begin{lrbox}{\SideBox}\begin{minipage}[t]{7.45cm}\raggedright\sloppy\color{Ink}\fontsize{12.2}{16.1}\selectfont \input{tikz/P04/D02.tex}\end{minipage}\end{lrbox}
\begin{lrbox}{\RecallBox}\begin{minipage}[t]{16.45cm}\raggedright\sloppy\color{Ink}\fontsize{10.1}{13.2}\selectfont La separazione riguarda tutte le decisioni che possono utilizzare l’informazione della risposta.\end{minipage}\end{lrbox}
\MeasureBox{5-TitleBox}{\TitleBox}
\MeasureBox{5-BodyBox}{\BodyBox}
\MeasureBox{5-SideBox}{\SideBox}
\MeasureBox{5-RecallBox}{\RecallBox}
\begin{frame}[plain]\hypertarget{slide5}{}
\bookmark[level=1,page=5]{Una separazione condivisa prima dell’esplorazione}
\begin{tikzpicture}[remember picture,overlay]\node[anchor=north east,inner sep=0] at ([xshift=19.05cm,yshift=-0.16cm]current page.north west){\color{Classical}\fontsize{8}{10}\selectfont\hyperlink{slide2}{Indice}};
\node[anchor=north west,inner sep=0] at ([xshift=0.95cm,yshift=-0.62cm]current page.north west){\usebox{\TitleBox}};
\draw[Classical,line width=1.3pt] ([xshift=.95cm,yshift=-1.68cm]current page.north west)--([xshift=2.25cm,yshift=-1.68cm]current page.north west);\node[anchor=north west,inner sep=0] at ([xshift=0.95cm,yshift=-2.12cm]current page.north west){\usebox{\BodyBox}};
\node[anchor=north west,inner sep=0] at ([xshift=11.6cm,yshift=-2.1cm]current page.north west){\color{Quantum}\bfseries\fontsize{11.5}{14}\selectfont Lettura del diagramma};
\node[anchor=north west,inner sep=0] at ([xshift=11.6cm,yshift=-2.75cm]current page.north west){\usebox{\SideBox}};
\draw[Rule,line width=.5pt] ([xshift=.95cm,yshift=-10.43cm]current page.north west)--([xshift=19.05cm,yshift=-10.43cm]current page.north west);\node[anchor=north west,inner sep=0] at ([xshift=0.95cm,yshift=-10.66cm]current page.north west){\color{Classical}\bfseries\fontsize{9.4}{12}\selectfont Richiami};
\node[anchor=north west,inner sep=0] at ([xshift=2.6cm,yshift=-10.64cm]current page.north west){\usebox{\RecallBox}};
\draw[Rule,line width=.35pt] ([xshift=.95cm,yshift=-11.92cm]current page.north west)--([xshift=19.05cm,yshift=-11.92cm]current page.north west);\node[anchor=north west,inner sep=0] at ([xshift=0.95cm,yshift=-12.06cm]current page.north west){\color{Muted}\fontsize{7.7}{9.3}\selectfont P04 · Un geyser: associazione e previsione};
\node[anchor=north east,inner sep=0] at ([xshift=17.65cm,yshift=-12.06cm]current page.north west){\color{Classical}\fontsize{7.3}{9.1}\selectfont };
\node[anchor=north east,inner sep=0] at ([xshift=19.05cm,yshift=-12.06cm]current page.north west){\color{Muted}\fontsize{7.7}{9.3}\selectfont 5/21};
\end{tikzpicture}\end{frame}
\begin{lrbox}{\TitleBox}\begin{minipage}[t]{18.2cm}\raggedright\sloppy\color{Ink}\fontsize{20.7}{24}\selectfont \bfseries\color{Navy}Identificativi e separazione verificabile\end{minipage}\end{lrbox}
\begin{lrbox}{\BodyBox}\begin{minipage}[t]{10.25cm}\raggedright\sloppy\color{Ink}\fontsize{12.4}{16.4}\selectfont La separazione viene applicata mediante l’identificativo, così che un ordinamento diverso della tabella non cambi l’appartenenza degli eventi. I controlli sul piano completo, eseguiti nello script integrale, verificano unicità e copertura di tutti gli identificativi prima dell’unione. Il frammento controlla poi numerosità e disgiunzione dei due insiemi. La riproduzione della procedura dipende dal piano salvato, mentre un semplice seme casuale condiviso fra linguaggi non garantirebbe la medesima assegnazione delle righe.\end{minipage}\end{lrbox}
\begin{lrbox}{\SideBox}\begin{minipage}[t]{7.45cm}\raggedright\sloppy\color{Ink}\fontsize{12.2}{16.1}\selectfont \par\vspace{.25cm}\lstinputlisting[language=R]{estratti/P04/separazione.R}\end{minipage}\end{lrbox}
\begin{lrbox}{\RecallBox}\begin{minipage}[t]{16.45cm}\raggedright\sloppy\color{Ink}\fontsize{10.1}{13.2}\selectfont Un piano esplicito conserva l’assegnazione degli eventi indipendentemente dall’ordinamento e dal generatore casuale.\end{minipage}\end{lrbox}
\MeasureBox{6-TitleBox}{\TitleBox}
\MeasureBox{6-BodyBox}{\BodyBox}
\MeasureBox{6-SideBox}{\SideBox}
\MeasureBox{6-RecallBox}{\RecallBox}
\begin{frame}[plain]\hypertarget{slide6}{}
\bookmark[level=1,page=6]{Identificativi e separazione verificabile}
\begin{tikzpicture}[remember picture,overlay]\node[anchor=north east,inner sep=0] at ([xshift=19.05cm,yshift=-0.16cm]current page.north west){\color{Classical}\fontsize{8}{10}\selectfont\hyperlink{slide2}{Indice}};
\node[anchor=north west,inner sep=0] at ([xshift=0.95cm,yshift=-0.62cm]current page.north west){\usebox{\TitleBox}};
\draw[Classical,line width=1.3pt] ([xshift=.95cm,yshift=-1.68cm]current page.north west)--([xshift=2.25cm,yshift=-1.68cm]current page.north west);\node[anchor=north west,inner sep=0] at ([xshift=0.95cm,yshift=-2.12cm]current page.north west){\usebox{\BodyBox}};
\node[anchor=north west,inner sep=0] at ([xshift=11.6cm,yshift=-2.1cm]current page.north west){\color{Quantum}\bfseries\fontsize{11.5}{14}\selectfont Implementazione in R};
\node[anchor=north west,inner sep=0] at ([xshift=11.6cm,yshift=-2.75cm]current page.north west){\usebox{\SideBox}};
\draw[Rule,line width=.5pt] ([xshift=.95cm,yshift=-10.43cm]current page.north west)--([xshift=19.05cm,yshift=-10.43cm]current page.north west);\node[anchor=north west,inner sep=0] at ([xshift=0.95cm,yshift=-10.66cm]current page.north west){\color{Classical}\bfseries\fontsize{9.4}{12}\selectfont Richiami};
\node[anchor=north west,inner sep=0] at ([xshift=2.6cm,yshift=-10.64cm]current page.north west){\usebox{\RecallBox}};
\draw[Rule,line width=.35pt] ([xshift=.95cm,yshift=-11.92cm]current page.north west)--([xshift=19.05cm,yshift=-11.92cm]current page.north west);\node[anchor=north west,inner sep=0] at ([xshift=0.95cm,yshift=-12.06cm]current page.north west){\color{Muted}\fontsize{7.7}{9.3}\selectfont P04 · Un geyser: associazione e previsione};
\node[anchor=north east,inner sep=0] at ([xshift=17.65cm,yshift=-12.06cm]current page.north west){\color{Classical}\fontsize{7.3}{9.1}\selectfont };
\node[anchor=north east,inner sep=0] at ([xshift=19.05cm,yshift=-12.06cm]current page.north west){\color{Muted}\fontsize{7.7}{9.3}\selectfont 6/21};
\end{tikzpicture}\end{frame}
\begin{lrbox}{\TitleBox}\begin{minipage}[t]{18.2cm}\raggedright\sloppy\color{Ink}\fontsize{20.7}{24}\selectfont \bfseries\color{Navy}La relazione mostra due addensamenti\end{minipage}\end{lrbox}
\begin{lrbox}{\BodyBox}\begin{minipage}[t]{10.25cm}\raggedright\sloppy\color{Ink}\fontsize{12.4}{16.4}\selectfont Il grafico delle 190 coppie di adattamento mostra una relazione crescente e due addensamenti nella durata e nell’attesa. Questi gruppi visivi suggeriscono una struttura che merita approfondimento, pur senza identificare automaticamente due stati fisici del sistema idrotermale. Una retta globale riassume insieme la separazione fra gli addensamenti e la variazione interna a ciascuno. La sua pendenza dovrà quindi essere interpretata come sintesi dell’associazione nel campione, verificando successivamente quanto differisca dalle pendenze entro intervalli di durata.\end{minipage}\end{lrbox}
\begin{lrbox}{\SideBox}\begin{minipage}[t]{7.45cm}\raggedright\sloppy\color{Ink}\fontsize{12.2}{16.1}\selectfont \includegraphics[width=7.45cm]{figure/P04/F01.pdf}\par\vspace{.16cm}{\fontsize{9.5}{12.6}\selectfont Coppie del solo insieme di adattamento. La scala dell’attesa segue la variabile pubblicata nella fonte.}\end{minipage}\end{lrbox}
\begin{lrbox}{\RecallBox}\begin{minipage}[t]{16.45cm}\raggedright\sloppy\color{Ink}\fontsize{10.1}{13.2}\selectfont La geometria del campione orienta l’interpretazione della pendenza prima della sua verifica inferenziale.\end{minipage}\end{lrbox}
\MeasureBox{7-TitleBox}{\TitleBox}
\MeasureBox{7-BodyBox}{\BodyBox}
\MeasureBox{7-SideBox}{\SideBox}
\MeasureBox{7-RecallBox}{\RecallBox}
\begin{frame}[plain]\hypertarget{slide7}{}
\bookmark[level=0,page=7]{Associazione nel campione di adattamento}
\bookmark[level=1,page=7]{La relazione mostra due addensamenti}
\begin{tikzpicture}[remember picture,overlay]\node[anchor=north east,inner sep=0] at ([xshift=19.05cm,yshift=-0.16cm]current page.north west){\color{Classical}\fontsize{8}{10}\selectfont\hyperlink{slide2}{Indice}};
\node[anchor=north west,inner sep=0] at ([xshift=0.95cm,yshift=-0.62cm]current page.north west){\usebox{\TitleBox}};
\draw[Classical,line width=1.3pt] ([xshift=.95cm,yshift=-1.68cm]current page.north west)--([xshift=2.25cm,yshift=-1.68cm]current page.north west);\node[anchor=north west,inner sep=0] at ([xshift=0.95cm,yshift=-2.12cm]current page.north west){\usebox{\BodyBox}};
\node[anchor=north west,inner sep=0] at ([xshift=11.6cm,yshift=-2.1cm]current page.north west){\color{Quantum}\bfseries\fontsize{11.5}{14}\selectfont Evidenza grafica};
\node[anchor=north west,inner sep=0] at ([xshift=11.6cm,yshift=-2.75cm]current page.north west){\usebox{\SideBox}};
\draw[Rule,line width=.5pt] ([xshift=.95cm,yshift=-10.43cm]current page.north west)--([xshift=19.05cm,yshift=-10.43cm]current page.north west);\node[anchor=north west,inner sep=0] at ([xshift=0.95cm,yshift=-10.66cm]current page.north west){\color{Classical}\bfseries\fontsize{9.4}{12}\selectfont Richiami};
\node[anchor=north west,inner sep=0] at ([xshift=2.6cm,yshift=-10.64cm]current page.north west){\usebox{\RecallBox}};
\draw[Rule,line width=.35pt] ([xshift=.95cm,yshift=-11.92cm]current page.north west)--([xshift=19.05cm,yshift=-11.92cm]current page.north west);\node[anchor=north west,inner sep=0] at ([xshift=0.95cm,yshift=-12.06cm]current page.north west){\color{Muted}\fontsize{7.7}{9.3}\selectfont P04 · Un geyser: associazione e previsione};
\node[anchor=north east,inner sep=0] at ([xshift=17.65cm,yshift=-12.06cm]current page.north west){\color{Classical}\fontsize{7.3}{9.1}\selectfont };
\node[anchor=north east,inner sep=0] at ([xshift=19.05cm,yshift=-12.06cm]current page.north west){\color{Muted}\fontsize{7.7}{9.3}\selectfont 7/21};
\end{tikzpicture}\end{frame}
\begin{lrbox}{\TitleBox}\begin{minipage}[t]{18.2cm}\raggedright\sloppy\color{Ink}\fontsize{20.7}{24}\selectfont \bfseries\color{Navy}La distribuzione delle durate completa il confronto\end{minipage}\end{lrbox}
\begin{lrbox}{\BodyBox}\begin{minipage}[t]{10.25cm}\raggedright\sloppy\color{Ink}\fontsize{12.4}{16.4}\selectfont L’istogramma rende esplicita la scarsa densità di eventi nella regione intermedia fra i due addensamenti. La posizione della retta in quella regione dipende dal modello lineare e riceve un sostegno osservativo minore rispetto alle zone più popolate. Gli intervalli delle barre sono fissati in quarti di minuto, mantenendo la scala e gli arrotondamenti della fonte. La lettura congiunta di istogramma e diagramma di dispersione permette di distinguere una relazione osservata diffusamente da una connessione sostenuta soprattutto da gruppi separati.\end{minipage}\end{lrbox}
\begin{lrbox}{\SideBox}\begin{minipage}[t]{7.45cm}\raggedright\sloppy\color{Ink}\fontsize{12.2}{16.1}\selectfont \includegraphics[width=7.45cm]{figure/P04/F02.pdf}\par\vspace{.16cm}{\fontsize{9.5}{12.6}\selectfont Distribuzione delle durate nei 190 eventi di adattamento; intervalli di ampiezza 0,25 minuti.}\end{minipage}\end{lrbox}
\begin{lrbox}{\RecallBox}\begin{minipage}[t]{16.45cm}\raggedright\sloppy\color{Ink}\fontsize{10.1}{13.2}\selectfont La densità dei predittori determina dove la relazione riceve maggiore sostegno empirico.\end{minipage}\end{lrbox}
\MeasureBox{8-TitleBox}{\TitleBox}
\MeasureBox{8-BodyBox}{\BodyBox}
\MeasureBox{8-SideBox}{\SideBox}
\MeasureBox{8-RecallBox}{\RecallBox}
\begin{frame}[plain]\hypertarget{slide8}{}
\bookmark[level=1,page=8]{La distribuzione delle durate completa il confronto}
\begin{tikzpicture}[remember picture,overlay]\node[anchor=north east,inner sep=0] at ([xshift=19.05cm,yshift=-0.16cm]current page.north west){\color{Classical}\fontsize{8}{10}\selectfont\hyperlink{slide2}{Indice}};
\node[anchor=north west,inner sep=0] at ([xshift=0.95cm,yshift=-0.62cm]current page.north west){\usebox{\TitleBox}};
\draw[Classical,line width=1.3pt] ([xshift=.95cm,yshift=-1.68cm]current page.north west)--([xshift=2.25cm,yshift=-1.68cm]current page.north west);\node[anchor=north west,inner sep=0] at ([xshift=0.95cm,yshift=-2.12cm]current page.north west){\usebox{\BodyBox}};
\node[anchor=north west,inner sep=0] at ([xshift=11.6cm,yshift=-2.1cm]current page.north west){\color{Quantum}\bfseries\fontsize{11.5}{14}\selectfont Evidenza grafica};
\node[anchor=north west,inner sep=0] at ([xshift=11.6cm,yshift=-2.75cm]current page.north west){\usebox{\SideBox}};
\draw[Rule,line width=.5pt] ([xshift=.95cm,yshift=-10.43cm]current page.north west)--([xshift=19.05cm,yshift=-10.43cm]current page.north west);\node[anchor=north west,inner sep=0] at ([xshift=0.95cm,yshift=-10.66cm]current page.north west){\color{Classical}\bfseries\fontsize{9.4}{12}\selectfont Richiami};
\node[anchor=north west,inner sep=0] at ([xshift=2.6cm,yshift=-10.64cm]current page.north west){\usebox{\RecallBox}};
\draw[Rule,line width=.35pt] ([xshift=.95cm,yshift=-11.92cm]current page.north west)--([xshift=19.05cm,yshift=-11.92cm]current page.north west);\node[anchor=north west,inner sep=0] at ([xshift=0.95cm,yshift=-12.06cm]current page.north west){\color{Muted}\fontsize{7.7}{9.3}\selectfont P04 · Un geyser: associazione e previsione};
\node[anchor=north east,inner sep=0] at ([xshift=17.65cm,yshift=-12.06cm]current page.north west){\color{Classical}\fontsize{7.3}{9.1}\selectfont };
\node[anchor=north east,inner sep=0] at ([xshift=19.05cm,yshift=-12.06cm]current page.north west){\color{Muted}\fontsize{7.7}{9.3}\selectfont 8/21};
\end{tikzpicture}\end{frame}
\begin{lrbox}{\TitleBox}\begin{minipage}[t]{18.2cm}\raggedright\sloppy\color{Ink}\fontsize{20.7}{24}\selectfont \bfseries\color{Navy}La pendenza misura un’associazione media\end{minipage}\end{lrbox}
\begin{lrbox}{\BodyBox}\begin{minipage}[t]{10.25cm}\raggedright\sloppy\color{Ink}\fontsize{12.4}{16.4}\selectfont La regressione collega l’attesa alla durata mediante una retta e un errore di media condizionata nulla. La pendenza stimata è circa 10,898 minuti di attesa per ogni minuto aggiuntivo di durata, mentre il coefficiente di determinazione è circa 0,819. Il coefficiente descrive l’associazione nella specificazione adottata e non identifica l’effetto di un intervento capace di prolungare un’eruzione. L’interpretazione fisica richiede un modello delle condizioni comuni che determinano le due variabili.\par\vspace{.28cm}{\centering\fontsize{13}{17}\selectfont\(\displaystyle\begin{gathered}Y_i=\alpha+\beta x_i+\varepsilon_i,\qquad E(\varepsilon_i\mid x_i)=0\\\widehat\beta=\frac{\sum_i(x_i-\overline x)(y_i-\overline y)}{\sum_i(x_i-\overline x)^2}\end{gathered}\)\par}\end{minipage}\end{lrbox}
\begin{lrbox}{\SideBox}\begin{minipage}[t]{7.45cm}\raggedright\sloppy\color{Ink}\fontsize{12.2}{16.1}\selectfont $x_i$: durata; $y_i$: attesa osservata, entrambe in minuti. Y indica la risposta casuale; $\alpha$: intercetta; $\beta$: pendenza; $\varepsilon_i$: errore. Le barre sono medie dei 190 eventi di adattamento. La pendenza ha unità min/min: il rapporto descrive due tempi distinti e non una variazione percentuale.\end{minipage}\end{lrbox}
\begin{lrbox}{\RecallBox}\begin{minipage}[t]{16.45cm}\raggedright\sloppy\color{Ink}\fontsize{10.1}{13.2}\selectfont L’interpretazione della pendenza conserva risposta, predittore e dominio del campione utilizzato.\end{minipage}\end{lrbox}
\MeasureBox{9-TitleBox}{\TitleBox}
\MeasureBox{9-BodyBox}{\BodyBox}
\MeasureBox{9-SideBox}{\SideBox}
\MeasureBox{9-RecallBox}{\RecallBox}
\begin{frame}[plain]\hypertarget{slide9}{}
\bookmark[level=1,page=9]{La pendenza misura un’associazione media}
\begin{tikzpicture}[remember picture,overlay]\node[anchor=north east,inner sep=0] at ([xshift=19.05cm,yshift=-0.16cm]current page.north west){\color{Classical}\fontsize{8}{10}\selectfont\hyperlink{slide2}{Indice}};
\node[anchor=north west,inner sep=0] at ([xshift=0.95cm,yshift=-0.62cm]current page.north west){\usebox{\TitleBox}};
\draw[Classical,line width=1.3pt] ([xshift=.95cm,yshift=-1.68cm]current page.north west)--([xshift=2.25cm,yshift=-1.68cm]current page.north west);\node[anchor=north west,inner sep=0] at ([xshift=0.95cm,yshift=-2.12cm]current page.north west){\usebox{\BodyBox}};
\node[anchor=north west,inner sep=0] at ([xshift=11.6cm,yshift=-2.1cm]current page.north west){\color{Quantum}\bfseries\fontsize{11.5}{14}\selectfont Sviluppo e interpretazione};
\node[anchor=north west,inner sep=0] at ([xshift=11.6cm,yshift=-2.75cm]current page.north west){\usebox{\SideBox}};
\draw[Rule,line width=.5pt] ([xshift=.95cm,yshift=-10.43cm]current page.north west)--([xshift=19.05cm,yshift=-10.43cm]current page.north west);\node[anchor=north west,inner sep=0] at ([xshift=0.95cm,yshift=-10.66cm]current page.north west){\color{Classical}\bfseries\fontsize{9.4}{12}\selectfont Richiami};
\node[anchor=north west,inner sep=0] at ([xshift=2.6cm,yshift=-10.64cm]current page.north west){\usebox{\RecallBox}};
\draw[Rule,line width=.35pt] ([xshift=.95cm,yshift=-11.92cm]current page.north west)--([xshift=19.05cm,yshift=-11.92cm]current page.north west);\node[anchor=north west,inner sep=0] at ([xshift=0.95cm,yshift=-12.06cm]current page.north west){\color{Muted}\fontsize{7.7}{9.3}\selectfont P04 · Un geyser: associazione e previsione};
\node[anchor=north east,inner sep=0] at ([xshift=17.65cm,yshift=-12.06cm]current page.north west){\color{Classical}\fontsize{7.3}{9.1}\selectfont };
\node[anchor=north east,inner sep=0] at ([xshift=19.05cm,yshift=-12.06cm]current page.north west){\color{Muted}\fontsize{7.7}{9.3}\selectfont 9/21};
\end{tikzpicture}\end{frame}
\begin{lrbox}{\TitleBox}\begin{minipage}[t]{18.2cm}\raggedright\sloppy\color{Ink}\fontsize{20.7}{24}\selectfont \bfseries\color{Navy}La covarianza HC3 considera varianze differenti\end{minipage}\end{lrbox}
\begin{lrbox}{\BodyBox}\begin{minipage}[t]{10.25cm}\raggedright\sloppy\color{Ink}\fontsize{12.4}{16.4}\selectfont La covarianza HC3 mantiene la stima OLS e pesa i residui attraverso la leva delle osservazioni. Il fattore uno meno la leva corregge la riduzione del residuo associata a punti che influenzano maggiormente il proprio adattamento. La formula ammette varianze diverse sotto ipotesi di indipendenza e regolarità. Non rappresenta invece automaticamente la dipendenza fra eventi consecutivi o gruppi non osservati, per i quali servirebbe una struttura di covarianza coerente con il disegno.\par\vspace{.28cm}{\centering\fontsize{13}{17}\selectfont\(\displaystyle\begin{gathered}h_i=v_i^T(X^TX)^{-1}v_i\\\widehat V_{HC3}=(X^TX)^{-1}X^T\operatorname{diag}\!\left(\frac{\widehat e_i^2}{(1-h_i)^2}\right)X(X^TX)^{-1}\end{gathered}\)\par}\end{minipage}\end{lrbox}
\begin{lrbox}{\SideBox}\begin{minipage}[t]{7.45cm}\raggedright\sloppy\color{Ink}\fontsize{12.2}{16.1}\selectfont \input{tikz/P04/D03.tex}\par\vspace{.16cm}{\fontsize{9.5}{12.6}\selectfont $v_i=(1,x_i)^T$: regressori; X ha righe $v_i^T$; $\hat e_i$: residuo; $h_i$: leva. HC3 corregge per eteroschedasticità.}\end{minipage}\end{lrbox}
\begin{lrbox}{\RecallBox}\begin{minipage}[t]{16.45cm}\raggedright\sloppy\color{Ink}\fontsize{10.1}{13.2}\selectfont Una covarianza robusta risponde a una forma dichiarata di deviazione dalle ipotesi classiche.\end{minipage}\end{lrbox}
\MeasureBox{10-TitleBox}{\TitleBox}
\MeasureBox{10-BodyBox}{\BodyBox}
\MeasureBox{10-SideBox}{\SideBox}
\MeasureBox{10-RecallBox}{\RecallBox}
\begin{frame}[plain]\hypertarget{slide10}{}
\bookmark[level=1,page=10]{La covarianza HC3 considera varianze differenti}
\begin{tikzpicture}[remember picture,overlay]\node[anchor=north east,inner sep=0] at ([xshift=19.05cm,yshift=-0.16cm]current page.north west){\color{Classical}\fontsize{8}{10}\selectfont\hyperlink{slide2}{Indice}};
\node[anchor=north west,inner sep=0] at ([xshift=0.95cm,yshift=-0.62cm]current page.north west){\usebox{\TitleBox}};
\draw[Classical,line width=1.3pt] ([xshift=.95cm,yshift=-1.68cm]current page.north west)--([xshift=2.25cm,yshift=-1.68cm]current page.north west);\node[anchor=north west,inner sep=0] at ([xshift=0.95cm,yshift=-2.12cm]current page.north west){\usebox{\BodyBox}};
\node[anchor=north west,inner sep=0] at ([xshift=11.6cm,yshift=-2.1cm]current page.north west){\color{Quantum}\bfseries\fontsize{11.5}{14}\selectfont Lettura del diagramma};
\node[anchor=north west,inner sep=0] at ([xshift=11.6cm,yshift=-2.75cm]current page.north west){\usebox{\SideBox}};
\draw[Rule,line width=.5pt] ([xshift=.95cm,yshift=-10.43cm]current page.north west)--([xshift=19.05cm,yshift=-10.43cm]current page.north west);\node[anchor=north west,inner sep=0] at ([xshift=0.95cm,yshift=-10.66cm]current page.north west){\color{Classical}\bfseries\fontsize{9.4}{12}\selectfont Richiami};
\node[anchor=north west,inner sep=0] at ([xshift=2.6cm,yshift=-10.64cm]current page.north west){\usebox{\RecallBox}};
\draw[Rule,line width=.35pt] ([xshift=.95cm,yshift=-11.92cm]current page.north west)--([xshift=19.05cm,yshift=-11.92cm]current page.north west);\node[anchor=north west,inner sep=0] at ([xshift=0.95cm,yshift=-12.06cm]current page.north west){\color{Muted}\fontsize{7.7}{9.3}\selectfont P04 · Un geyser: associazione e previsione};
\node[anchor=north east,inner sep=0] at ([xshift=17.65cm,yshift=-12.06cm]current page.north west){\color{Classical}\fontsize{7.3}{9.1}\selectfont MacKinnon–White; statsmodels, HC3\_se};
\node[anchor=north east,inner sep=0] at ([xshift=19.05cm,yshift=-12.06cm]current page.north west){\color{Muted}\fontsize{7.7}{9.3}\selectfont 10/21};
\end{tikzpicture}\end{frame}
\begin{lrbox}{\TitleBox}\begin{minipage}[t]{18.2cm}\raggedright\sloppy\color{Ink}\fontsize{20.7}{24}\selectfont \bfseries\color{Navy}Confrontare gli intervalli della pendenza\end{minipage}\end{lrbox}
\begin{lrbox}{\BodyBox}\begin{minipage}[t]{10.25cm}\raggedright\sloppy\color{Ink}\fontsize{12.4}{16.4}\selectfont L’intervallo t classico al 95\% va da circa 10,161 a 11,635, mentre la costruzione HC3 va da circa 10,188 a 11,608 minuti per minuto. Nel secondo caso si utilizza lo stesso quantile t con 188 gradi di libertà come approssimazione dichiarata. Le ampiezze risultano simili in questo campione, senza che ciò costituisca una prova di omoschedasticità. La vicinanza fra le due costruzioni documenta soltanto una sensibilità limitata alla variante di covarianza effettivamente esaminata.\end{minipage}\end{lrbox}
\begin{lrbox}{\SideBox}\begin{minipage}[t]{7.45cm}\raggedright\sloppy\color{Ink}\fontsize{12.2}{16.1}\selectfont \includegraphics[width=7.45cm]{figure/P04/F09.pdf}\par\vspace{.16cm}{\fontsize{9.5}{12.6}\selectfont Stima OLS comune e due intervalli della pendenza. Per HC3 il riferimento t è approssimativo.}\end{minipage}\end{lrbox}
\begin{lrbox}{\RecallBox}\begin{minipage}[t]{16.45cm}\raggedright\sloppy\color{Ink}\fontsize{10.1}{13.2}\selectfont La concordanza fra due procedure viene interpretata entro le ipotesi e le varianti considerate.\end{minipage}\end{lrbox}
\MeasureBox{11-TitleBox}{\TitleBox}
\MeasureBox{11-BodyBox}{\BodyBox}
\MeasureBox{11-SideBox}{\SideBox}
\MeasureBox{11-RecallBox}{\RecallBox}
\begin{frame}[plain]\hypertarget{slide11}{}
\bookmark[level=1,page=11]{Confrontare gli intervalli della pendenza}
\begin{tikzpicture}[remember picture,overlay]\node[anchor=north east,inner sep=0] at ([xshift=19.05cm,yshift=-0.16cm]current page.north west){\color{Classical}\fontsize{8}{10}\selectfont\hyperlink{slide2}{Indice}};
\node[anchor=north west,inner sep=0] at ([xshift=0.95cm,yshift=-0.62cm]current page.north west){\usebox{\TitleBox}};
\draw[Classical,line width=1.3pt] ([xshift=.95cm,yshift=-1.68cm]current page.north west)--([xshift=2.25cm,yshift=-1.68cm]current page.north west);\node[anchor=north west,inner sep=0] at ([xshift=0.95cm,yshift=-2.12cm]current page.north west){\usebox{\BodyBox}};
\node[anchor=north west,inner sep=0] at ([xshift=11.6cm,yshift=-2.1cm]current page.north west){\color{Quantum}\bfseries\fontsize{11.5}{14}\selectfont Evidenza grafica};
\node[anchor=north west,inner sep=0] at ([xshift=11.6cm,yshift=-2.75cm]current page.north west){\usebox{\SideBox}};
\draw[Rule,line width=.5pt] ([xshift=.95cm,yshift=-10.43cm]current page.north west)--([xshift=19.05cm,yshift=-10.43cm]current page.north west);\node[anchor=north west,inner sep=0] at ([xshift=0.95cm,yshift=-10.66cm]current page.north west){\color{Classical}\bfseries\fontsize{9.4}{12}\selectfont Richiami};
\node[anchor=north west,inner sep=0] at ([xshift=2.6cm,yshift=-10.64cm]current page.north west){\usebox{\RecallBox}};
\draw[Rule,line width=.35pt] ([xshift=.95cm,yshift=-11.92cm]current page.north west)--([xshift=19.05cm,yshift=-11.92cm]current page.north west);\node[anchor=north west,inner sep=0] at ([xshift=0.95cm,yshift=-12.06cm]current page.north west){\color{Muted}\fontsize{7.7}{9.3}\selectfont P04 · Un geyser: associazione e previsione};
\node[anchor=north east,inner sep=0] at ([xshift=17.65cm,yshift=-12.06cm]current page.north west){\color{Classical}\fontsize{7.3}{9.1}\selectfont };
\node[anchor=north east,inner sep=0] at ([xshift=19.05cm,yshift=-12.06cm]current page.north west){\color{Muted}\fontsize{7.7}{9.3}\selectfont 11/21};
\end{tikzpicture}\end{frame}
\begin{lrbox}{\TitleBox}\begin{minipage}[t]{18.2cm}\raggedright\sloppy\color{Ink}\fontsize{20.7}{24}\selectfont \bfseries\color{Navy}L’incertezza di un nuovo evento\end{minipage}\end{lrbox}
\begin{lrbox}{\BodyBox}\begin{minipage}[t]{10.25cm}\raggedright\sloppy\color{Ink}\fontsize{12.4}{16.4}\selectfont L’intervallo della media descrive l’incertezza sulla posizione della retta a una durata fissata. Per un nuovo evento si aggiunge la variabilità residua individuale, che compare nella formula attraverso il termine uno sotto radice. Entrambe le costruzioni adottano il modello classico omoschedastico e la stessa varianza residua. La covarianza HC3 della pendenza non viene trasformata automaticamente in un intervallo predittivo robusto, poiché quest’ultimo richiederebbe anche un modello della distribuzione del nuovo errore.\par\vspace{.28cm}{\centering\fontsize{13}{17}\selectfont\(\displaystyle\begin{gathered}h_0=v_0^T(X^TX)^{-1}v_0\\IC_{\rm media}:\ \hat y_0\pm ts\sqrt{h_0},\quad IP:\ \hat y_0\pm ts\sqrt{1+h_0}\end{gathered}\)\par}\end{minipage}\end{lrbox}
\begin{lrbox}{\SideBox}\begin{minipage}[t]{7.45cm}\raggedright\sloppy\color{Ink}\fontsize{12.2}{16.1}\selectfont \input{tikz/P04/D04.tex}\par\vspace{.16cm}{\fontsize{9.5}{12.6}\selectfont $v_0=(1,x_0)^T$; X: colonne 1 e durata; $\hat y_0$: attesa stimata. $s^2=\sum\hat e_i^2/188$; $t=t_{0,975;188}$.}\end{minipage}\end{lrbox}
\begin{lrbox}{\RecallBox}\begin{minipage}[t]{16.45cm}\raggedright\sloppy\color{Ink}\fontsize{10.1}{13.2}\selectfont La quantità oggetto dell’intervallo deve essere definita prima di interpretarne l’ampiezza.\end{minipage}\end{lrbox}
\MeasureBox{12-TitleBox}{\TitleBox}
\MeasureBox{12-BodyBox}{\BodyBox}
\MeasureBox{12-SideBox}{\SideBox}
\MeasureBox{12-RecallBox}{\RecallBox}
\begin{frame}[plain]\hypertarget{slide12}{}
\bookmark[level=0,page=12]{Incertezza della media e di un evento}
\bookmark[level=1,page=12]{L’incertezza di un nuovo evento}
\begin{tikzpicture}[remember picture,overlay]\node[anchor=north east,inner sep=0] at ([xshift=19.05cm,yshift=-0.16cm]current page.north west){\color{Classical}\fontsize{8}{10}\selectfont\hyperlink{slide2}{Indice}};
\node[anchor=north west,inner sep=0] at ([xshift=0.95cm,yshift=-0.62cm]current page.north west){\usebox{\TitleBox}};
\draw[Classical,line width=1.3pt] ([xshift=.95cm,yshift=-1.68cm]current page.north west)--([xshift=2.25cm,yshift=-1.68cm]current page.north west);\node[anchor=north west,inner sep=0] at ([xshift=0.95cm,yshift=-2.12cm]current page.north west){\usebox{\BodyBox}};
\node[anchor=north west,inner sep=0] at ([xshift=11.6cm,yshift=-2.1cm]current page.north west){\color{Quantum}\bfseries\fontsize{11.5}{14}\selectfont Lettura del diagramma};
\node[anchor=north west,inner sep=0] at ([xshift=11.6cm,yshift=-2.75cm]current page.north west){\usebox{\SideBox}};
\draw[Rule,line width=.5pt] ([xshift=.95cm,yshift=-10.43cm]current page.north west)--([xshift=19.05cm,yshift=-10.43cm]current page.north west);\node[anchor=north west,inner sep=0] at ([xshift=0.95cm,yshift=-10.66cm]current page.north west){\color{Classical}\bfseries\fontsize{9.4}{12}\selectfont Richiami};
\node[anchor=north west,inner sep=0] at ([xshift=2.6cm,yshift=-10.64cm]current page.north west){\usebox{\RecallBox}};
\draw[Rule,line width=.35pt] ([xshift=.95cm,yshift=-11.92cm]current page.north west)--([xshift=19.05cm,yshift=-11.92cm]current page.north west);\node[anchor=north west,inner sep=0] at ([xshift=0.95cm,yshift=-12.06cm]current page.north west){\color{Muted}\fontsize{7.7}{9.3}\selectfont P04 · Un geyser: associazione e previsione};
\node[anchor=north east,inner sep=0] at ([xshift=17.65cm,yshift=-12.06cm]current page.north west){\color{Classical}\fontsize{7.3}{9.1}\selectfont };
\node[anchor=north east,inner sep=0] at ([xshift=19.05cm,yshift=-12.06cm]current page.north west){\color{Muted}\fontsize{7.7}{9.3}\selectfont 12/21};
\end{tikzpicture}\end{frame}
\begin{lrbox}{\TitleBox}\begin{minipage}[t]{18.2cm}\raggedright\sloppy\color{Ink}\fontsize{20.7}{24}\selectfont \bfseries\color{Navy}Le due fasce vengono lette sulla stessa scala\end{minipage}\end{lrbox}
\begin{lrbox}{\BodyBox}\begin{minipage}[t]{10.25cm}\raggedright\sloppy\color{Ink}\fontsize{12.4}{16.4}\selectfont La fascia verde segue l’incertezza della media condizionata, mentre quella viola rappresenta il nuovo evento sotto le ipotesi classiche. Entrambe si allargano allontanandosi dalla durata media del campione, poiché cresce la leva del punto di previsione. La fascia individuale rimane più ampia per il contributo della variabilità residua. Il disegno viene limitato al dominio delle durate osservate nel campione di adattamento, entro il quale restano comunque zone poco popolate e assunzioni di forma che devono essere discusse.\end{minipage}\end{lrbox}
\begin{lrbox}{\SideBox}\begin{minipage}[t]{7.45cm}\raggedright\sloppy\color{Ink}\fontsize{12.2}{16.1}\selectfont \includegraphics[width=7.45cm]{figure/P04/F03.pdf}\par\vspace{.16cm}{\fontsize{9.5}{12.6}\selectfont Verde: intervallo della media; viola: intervallo predittivo individuale. Costruzioni classiche al 95\%.}\end{minipage}\end{lrbox}
\begin{lrbox}{\RecallBox}\begin{minipage}[t]{16.45cm}\raggedright\sloppy\color{Ink}\fontsize{10.1}{13.2}\selectfont Una fascia stretta intorno alla retta non implica una previsione individuale altrettanto precisa.\end{minipage}\end{lrbox}
\MeasureBox{13-TitleBox}{\TitleBox}
\MeasureBox{13-BodyBox}{\BodyBox}
\MeasureBox{13-SideBox}{\SideBox}
\MeasureBox{13-RecallBox}{\RecallBox}
\begin{frame}[plain]\hypertarget{slide13}{}
\bookmark[level=1,page=13]{Le due fasce vengono lette sulla stessa scala}
\begin{tikzpicture}[remember picture,overlay]\node[anchor=north east,inner sep=0] at ([xshift=19.05cm,yshift=-0.16cm]current page.north west){\color{Classical}\fontsize{8}{10}\selectfont\hyperlink{slide2}{Indice}};
\node[anchor=north west,inner sep=0] at ([xshift=0.95cm,yshift=-0.62cm]current page.north west){\usebox{\TitleBox}};
\draw[Classical,line width=1.3pt] ([xshift=.95cm,yshift=-1.68cm]current page.north west)--([xshift=2.25cm,yshift=-1.68cm]current page.north west);\node[anchor=north west,inner sep=0] at ([xshift=0.95cm,yshift=-2.12cm]current page.north west){\usebox{\BodyBox}};
\node[anchor=north west,inner sep=0] at ([xshift=11.6cm,yshift=-2.1cm]current page.north west){\color{Quantum}\bfseries\fontsize{11.5}{14}\selectfont Evidenza grafica};
\node[anchor=north west,inner sep=0] at ([xshift=11.6cm,yshift=-2.75cm]current page.north west){\usebox{\SideBox}};
\draw[Rule,line width=.5pt] ([xshift=.95cm,yshift=-10.43cm]current page.north west)--([xshift=19.05cm,yshift=-10.43cm]current page.north west);\node[anchor=north west,inner sep=0] at ([xshift=0.95cm,yshift=-10.66cm]current page.north west){\color{Classical}\bfseries\fontsize{9.4}{12}\selectfont Richiami};
\node[anchor=north west,inner sep=0] at ([xshift=2.6cm,yshift=-10.64cm]current page.north west){\usebox{\RecallBox}};
\draw[Rule,line width=.35pt] ([xshift=.95cm,yshift=-11.92cm]current page.north west)--([xshift=19.05cm,yshift=-11.92cm]current page.north west);\node[anchor=north west,inner sep=0] at ([xshift=0.95cm,yshift=-12.06cm]current page.north west){\color{Muted}\fontsize{7.7}{9.3}\selectfont P04 · Un geyser: associazione e previsione};
\node[anchor=north east,inner sep=0] at ([xshift=17.65cm,yshift=-12.06cm]current page.north west){\color{Classical}\fontsize{7.3}{9.1}\selectfont };
\node[anchor=north east,inner sep=0] at ([xshift=19.05cm,yshift=-12.06cm]current page.north west){\color{Muted}\fontsize{7.7}{9.3}\selectfont 13/21};
\end{tikzpicture}\end{frame}
\begin{lrbox}{\TitleBox}\begin{minipage}[t]{18.2cm}\raggedright\sloppy\color{Ink}\fontsize{20.7}{24}\selectfont \bfseries\color{Navy}I residui precedono l’interpretazione predittiva\end{minipage}\end{lrbox}
\begin{lrbox}{\BodyBox}\begin{minipage}[t]{10.25cm}\raggedright\sloppy\color{Ink}\fontsize{12.4}{16.4}\selectfont Il grafico dei residui rispetto alle attese stimate esamina la forma della relazione e la dispersione rimasta attorno alla retta. Strutture sistematiche potrebbero suggerire curvature, gruppi o variazioni di varianza, mentre punti isolati richiederebbero un controllo della fonte. La diagnosi resta confinata ai dati di adattamento. Le proprietà algebriche dei residui OLS, fra cui l’ortogonalità ai predittori, non garantiscono da sole la forma distributiva dell’errore futuro né la stabilità della relazione fra campagne osservative differenti.\end{minipage}\end{lrbox}
\begin{lrbox}{\SideBox}\begin{minipage}[t]{7.45cm}\raggedright\sloppy\color{Ink}\fontsize{12.2}{16.1}\selectfont \includegraphics[width=7.45cm]{figure/P04/F04.pdf}\par\vspace{.16cm}{\fontsize{9.5}{12.6}\selectfont Residui dei 190 eventi usati per stimare la retta. Il riferimento tratteggiato indica residuo nullo.}\end{minipage}\end{lrbox}
\begin{lrbox}{\RecallBox}\begin{minipage}[t]{16.45cm}\raggedright\sloppy\color{Ink}\fontsize{10.1}{13.2}\selectfont La diagnosi del modello interpreta la struttura residua prima della valutazione sui dati separati.\end{minipage}\end{lrbox}
\MeasureBox{14-TitleBox}{\TitleBox}
\MeasureBox{14-BodyBox}{\BodyBox}
\MeasureBox{14-SideBox}{\SideBox}
\MeasureBox{14-RecallBox}{\RecallBox}
\begin{frame}[plain]\hypertarget{slide14}{}
\bookmark[level=1,page=14]{I residui precedono l’interpretazione predittiva}
\begin{tikzpicture}[remember picture,overlay]\node[anchor=north east,inner sep=0] at ([xshift=19.05cm,yshift=-0.16cm]current page.north west){\color{Classical}\fontsize{8}{10}\selectfont\hyperlink{slide2}{Indice}};
\node[anchor=north west,inner sep=0] at ([xshift=0.95cm,yshift=-0.62cm]current page.north west){\usebox{\TitleBox}};
\draw[Classical,line width=1.3pt] ([xshift=.95cm,yshift=-1.68cm]current page.north west)--([xshift=2.25cm,yshift=-1.68cm]current page.north west);\node[anchor=north west,inner sep=0] at ([xshift=0.95cm,yshift=-2.12cm]current page.north west){\usebox{\BodyBox}};
\node[anchor=north west,inner sep=0] at ([xshift=11.6cm,yshift=-2.1cm]current page.north west){\color{Quantum}\bfseries\fontsize{11.5}{14}\selectfont Evidenza grafica};
\node[anchor=north west,inner sep=0] at ([xshift=11.6cm,yshift=-2.75cm]current page.north west){\usebox{\SideBox}};
\draw[Rule,line width=.5pt] ([xshift=.95cm,yshift=-10.43cm]current page.north west)--([xshift=19.05cm,yshift=-10.43cm]current page.north west);\node[anchor=north west,inner sep=0] at ([xshift=0.95cm,yshift=-10.66cm]current page.north west){\color{Classical}\bfseries\fontsize{9.4}{12}\selectfont Richiami};
\node[anchor=north west,inner sep=0] at ([xshift=2.6cm,yshift=-10.64cm]current page.north west){\usebox{\RecallBox}};
\draw[Rule,line width=.35pt] ([xshift=.95cm,yshift=-11.92cm]current page.north west)--([xshift=19.05cm,yshift=-11.92cm]current page.north west);\node[anchor=north west,inner sep=0] at ([xshift=0.95cm,yshift=-12.06cm]current page.north west){\color{Muted}\fontsize{7.7}{9.3}\selectfont P04 · Un geyser: associazione e previsione};
\node[anchor=north east,inner sep=0] at ([xshift=17.65cm,yshift=-12.06cm]current page.north west){\color{Classical}\fontsize{7.3}{9.1}\selectfont };
\node[anchor=north east,inner sep=0] at ([xshift=19.05cm,yshift=-12.06cm]current page.north west){\color{Muted}\fontsize{7.7}{9.3}\selectfont 14/21};
\end{tikzpicture}\end{frame}
\begin{lrbox}{\TitleBox}\begin{minipage}[t]{18.2cm}\raggedright\sloppy\color{Ink}\fontsize{20.7}{24}\selectfont \bfseries\color{Navy}La previsione costante usa soltanto l’adattamento\end{minipage}\end{lrbox}
\begin{lrbox}{\BodyBox}\begin{minipage}[t]{10.25cm}\raggedright\sloppy\color{Ink}\fontsize{12.4}{16.4}\selectfont La procedura di riferimento assegna a ogni evento di verifica la media dell’attesa calcolata sui 190 eventi di adattamento, circa 70,479 minuti. Questa scelta stabilisce quanto si guadagni utilizzando la durata rispetto a una previsione priva di predittori. Impiegare la media degli 82 eventi di verifica introdurrebbe informazione della risposta proprio nel riferimento che si intende valutare. Entrambe le procedure vengono quindi definite prima del confronto finale e misurate sulle stesse osservazioni.\par\vspace{.28cm}{\centering\fontsize{13}{17}\selectfont\(\displaystyle\begin{gathered}\widehat y_i^{(0)}=\overline y_{train},\qquad \widehat y_i^{(1)}=\widehat\alpha+\widehat\beta x_i\end{gathered}\)\par}\end{minipage}\end{lrbox}
\begin{lrbox}{\SideBox}\begin{minipage}[t]{7.45cm}\raggedright\sloppy\color{Ink}\fontsize{12.2}{16.1}\selectfont \input{tikz/P04/D05.tex}\par\vspace{.16cm}{\fontsize{9.5}{12.6}\selectfont $x_i$: durata del nuovo evento; $\bar y_{train}$: attesa media; $\hat\alpha,\hat\beta$: coefficienti. Tutte le stime usano i soli 190 eventi di adattamento.}\end{minipage}\end{lrbox}
\begin{lrbox}{\RecallBox}\begin{minipage}[t]{16.45cm}\raggedright\sloppy\color{Ink}\fontsize{10.1}{13.2}\selectfont Il riferimento predittivo deve rispettare gli stessi vincoli informativi del modello esaminato.\end{minipage}\end{lrbox}
\MeasureBox{15-TitleBox}{\TitleBox}
\MeasureBox{15-BodyBox}{\BodyBox}
\MeasureBox{15-SideBox}{\SideBox}
\MeasureBox{15-RecallBox}{\RecallBox}
\begin{frame}[plain]\hypertarget{slide15}{}
\bookmark[level=0,page=15]{Valutazione su eventi separati}
\bookmark[level=1,page=15]{La previsione costante usa soltanto l’adattamento}
\begin{tikzpicture}[remember picture,overlay]\node[anchor=north east,inner sep=0] at ([xshift=19.05cm,yshift=-0.16cm]current page.north west){\color{Classical}\fontsize{8}{10}\selectfont\hyperlink{slide2}{Indice}};
\node[anchor=north west,inner sep=0] at ([xshift=0.95cm,yshift=-0.62cm]current page.north west){\usebox{\TitleBox}};
\draw[Classical,line width=1.3pt] ([xshift=.95cm,yshift=-1.68cm]current page.north west)--([xshift=2.25cm,yshift=-1.68cm]current page.north west);\node[anchor=north west,inner sep=0] at ([xshift=0.95cm,yshift=-2.12cm]current page.north west){\usebox{\BodyBox}};
\node[anchor=north west,inner sep=0] at ([xshift=11.6cm,yshift=-2.1cm]current page.north west){\color{Quantum}\bfseries\fontsize{11.5}{14}\selectfont Lettura del diagramma};
\node[anchor=north west,inner sep=0] at ([xshift=11.6cm,yshift=-2.75cm]current page.north west){\usebox{\SideBox}};
\draw[Rule,line width=.5pt] ([xshift=.95cm,yshift=-10.43cm]current page.north west)--([xshift=19.05cm,yshift=-10.43cm]current page.north west);\node[anchor=north west,inner sep=0] at ([xshift=0.95cm,yshift=-10.66cm]current page.north west){\color{Classical}\bfseries\fontsize{9.4}{12}\selectfont Richiami};
\node[anchor=north west,inner sep=0] at ([xshift=2.6cm,yshift=-10.64cm]current page.north west){\usebox{\RecallBox}};
\draw[Rule,line width=.35pt] ([xshift=.95cm,yshift=-11.92cm]current page.north west)--([xshift=19.05cm,yshift=-11.92cm]current page.north west);\node[anchor=north west,inner sep=0] at ([xshift=0.95cm,yshift=-12.06cm]current page.north west){\color{Muted}\fontsize{7.7}{9.3}\selectfont P04 · Un geyser: associazione e previsione};
\node[anchor=north east,inner sep=0] at ([xshift=17.65cm,yshift=-12.06cm]current page.north west){\color{Classical}\fontsize{7.3}{9.1}\selectfont };
\node[anchor=north east,inner sep=0] at ([xshift=19.05cm,yshift=-12.06cm]current page.north west){\color{Muted}\fontsize{7.7}{9.3}\selectfont 15/21};
\end{tikzpicture}\end{frame}
\begin{lrbox}{\TitleBox}\begin{minipage}[t]{18.2cm}\raggedright\sloppy\color{Ink}\fontsize{20.7}{24}\selectfont \bfseries\color{Navy}La misura dell’errore mantiene il campione separato\end{minipage}\end{lrbox}
\begin{lrbox}{\BodyBox}\begin{minipage}[t]{10.25cm}\raggedright\sloppy\color{Ink}\fontsize{12.4}{16.4}\selectfont Il codice adatta la retta sui dati train e produce previsioni per le durate presenti nel test. Lo scarto fra attesa osservata e prevista viene poi riassunto mediante RMSE e MAE, che attribuiscono peso diverso agli errori grandi. Le misure riguardano gli 82 eventi separati. Una modifica del modello suggerita dai risultati del test richiederebbe un ulteriore campione indipendente per una nuova valutazione finale.\par\vspace{.28cm}{\centering\fontsize{13}{17}\selectfont\(\displaystyle\begin{gathered}RMSE=\sqrt{\frac1{82}\sum_i(y_i-\widehat y_i)^2}\\MAE=\frac1{82}\sum_i|y_i-\widehat y_i|\end{gathered}\)\par}\end{minipage}\end{lrbox}
\begin{lrbox}{\SideBox}\begin{minipage}[t]{7.45cm}\raggedright\sloppy\color{Ink}\fontsize{12.2}{16.1}\selectfont $y_i$: attesa osservata; $\hat y_i$: previsione, per i = 1,…,82 nel test. RMSE e MAE sono entrambe in minuti.\par\vspace{.25cm}\lstinputlisting[language=R]{estratti/P04/previsione.R}\end{minipage}\end{lrbox}
\begin{lrbox}{\RecallBox}\begin{minipage}[t]{16.45cm}\raggedright\sloppy\color{Ink}\fontsize{10.1}{13.2}\selectfont Il test valuta una procedura definita; il suo riutilizzo per scegliere la procedura modifica il disegno.\end{minipage}\end{lrbox}
\MeasureBox{16-TitleBox}{\TitleBox}
\MeasureBox{16-BodyBox}{\BodyBox}
\MeasureBox{16-SideBox}{\SideBox}
\MeasureBox{16-RecallBox}{\RecallBox}
\begin{frame}[plain]\hypertarget{slide16}{}
\bookmark[level=1,page=16]{La misura dell’errore mantiene il campione separato}
\begin{tikzpicture}[remember picture,overlay]\node[anchor=north east,inner sep=0] at ([xshift=19.05cm,yshift=-0.16cm]current page.north west){\color{Classical}\fontsize{8}{10}\selectfont\hyperlink{slide2}{Indice}};
\node[anchor=north west,inner sep=0] at ([xshift=0.95cm,yshift=-0.62cm]current page.north west){\usebox{\TitleBox}};
\draw[Classical,line width=1.3pt] ([xshift=.95cm,yshift=-1.68cm]current page.north west)--([xshift=2.25cm,yshift=-1.68cm]current page.north west);\node[anchor=north west,inner sep=0] at ([xshift=0.95cm,yshift=-2.12cm]current page.north west){\usebox{\BodyBox}};
\node[anchor=north west,inner sep=0] at ([xshift=11.6cm,yshift=-2.1cm]current page.north west){\color{Quantum}\bfseries\fontsize{11.5}{14}\selectfont Implementazione in R};
\node[anchor=north west,inner sep=0] at ([xshift=11.6cm,yshift=-2.75cm]current page.north west){\usebox{\SideBox}};
\draw[Rule,line width=.5pt] ([xshift=.95cm,yshift=-10.43cm]current page.north west)--([xshift=19.05cm,yshift=-10.43cm]current page.north west);\node[anchor=north west,inner sep=0] at ([xshift=0.95cm,yshift=-10.66cm]current page.north west){\color{Classical}\bfseries\fontsize{9.4}{12}\selectfont Richiami};
\node[anchor=north west,inner sep=0] at ([xshift=2.6cm,yshift=-10.64cm]current page.north west){\usebox{\RecallBox}};
\draw[Rule,line width=.35pt] ([xshift=.95cm,yshift=-11.92cm]current page.north west)--([xshift=19.05cm,yshift=-11.92cm]current page.north west);\node[anchor=north west,inner sep=0] at ([xshift=0.95cm,yshift=-12.06cm]current page.north west){\color{Muted}\fontsize{7.7}{9.3}\selectfont P04 · Un geyser: associazione e previsione};
\node[anchor=north east,inner sep=0] at ([xshift=17.65cm,yshift=-12.06cm]current page.north west){\color{Classical}\fontsize{7.3}{9.1}\selectfont };
\node[anchor=north east,inner sep=0] at ([xshift=19.05cm,yshift=-12.06cm]current page.north west){\color{Muted}\fontsize{7.7}{9.3}\selectfont 16/21};
\end{tikzpicture}\end{frame}
\begin{lrbox}{\TitleBox}\begin{minipage}[t]{18.2cm}\raggedright\sloppy\color{Ink}\fontsize{20.7}{24}\selectfont \bfseries\color{Navy}Osservato e previsto rendono visibili gli scarti\end{minipage}\end{lrbox}
\begin{lrbox}{\BodyBox}\begin{minipage}[t]{10.25cm}\raggedright\sloppy\color{Ink}\fontsize{12.4}{16.4}\selectfont Il diagramma confronta ogni attesa osservata con la previsione corrispondente, adottando la stessa scala sui due assi. La bisettrice rappresenta previsione esatta, mentre la distanza verticale rende leggibili sovrastime e sottostime. La presenza dei due addensamenti resta visibile anche negli eventi separati. Una correlazione elevata fra osservato e previsto sarebbe insufficiente a descrivere accuratezza e calibrazione, poiché potrebbe convivere con errori sistematici; per questo il confronto viene accompagnato da misure nella scala dei minuti.\end{minipage}\end{lrbox}
\begin{lrbox}{\SideBox}\begin{minipage}[t]{7.45cm}\raggedright\sloppy\color{Ink}\fontsize{12.2}{16.1}\selectfont \includegraphics[width=7.45cm]{figure/P04/F05.pdf}\par\vspace{.16cm}{\fontsize{9.5}{12.6}\selectfont Ottantadue eventi di verifica. Bisettrice: previsione uguale all’osservazione; rapporto unitario fra gli assi.}\end{minipage}\end{lrbox}
\begin{lrbox}{\RecallBox}\begin{minipage}[t]{16.45cm}\raggedright\sloppy\color{Ink}\fontsize{10.1}{13.2}\selectfont La valutazione predittiva richiede errori espressi nella scala della risposta, oltre all’associazione grafica.\end{minipage}\end{lrbox}
\MeasureBox{17-TitleBox}{\TitleBox}
\MeasureBox{17-BodyBox}{\BodyBox}
\MeasureBox{17-SideBox}{\SideBox}
\MeasureBox{17-RecallBox}{\RecallBox}
\begin{frame}[plain]\hypertarget{slide17}{}
\bookmark[level=1,page=17]{Osservato e previsto rendono visibili gli scarti}
\begin{tikzpicture}[remember picture,overlay]\node[anchor=north east,inner sep=0] at ([xshift=19.05cm,yshift=-0.16cm]current page.north west){\color{Classical}\fontsize{8}{10}\selectfont\hyperlink{slide2}{Indice}};
\node[anchor=north west,inner sep=0] at ([xshift=0.95cm,yshift=-0.62cm]current page.north west){\usebox{\TitleBox}};
\draw[Classical,line width=1.3pt] ([xshift=.95cm,yshift=-1.68cm]current page.north west)--([xshift=2.25cm,yshift=-1.68cm]current page.north west);\node[anchor=north west,inner sep=0] at ([xshift=0.95cm,yshift=-2.12cm]current page.north west){\usebox{\BodyBox}};
\node[anchor=north west,inner sep=0] at ([xshift=11.6cm,yshift=-2.1cm]current page.north west){\color{Quantum}\bfseries\fontsize{11.5}{14}\selectfont Evidenza grafica};
\node[anchor=north west,inner sep=0] at ([xshift=11.6cm,yshift=-2.75cm]current page.north west){\usebox{\SideBox}};
\draw[Rule,line width=.5pt] ([xshift=.95cm,yshift=-10.43cm]current page.north west)--([xshift=19.05cm,yshift=-10.43cm]current page.north west);\node[anchor=north west,inner sep=0] at ([xshift=0.95cm,yshift=-10.66cm]current page.north west){\color{Classical}\bfseries\fontsize{9.4}{12}\selectfont Richiami};
\node[anchor=north west,inner sep=0] at ([xshift=2.6cm,yshift=-10.64cm]current page.north west){\usebox{\RecallBox}};
\draw[Rule,line width=.35pt] ([xshift=.95cm,yshift=-11.92cm]current page.north west)--([xshift=19.05cm,yshift=-11.92cm]current page.north west);\node[anchor=north west,inner sep=0] at ([xshift=0.95cm,yshift=-12.06cm]current page.north west){\color{Muted}\fontsize{7.7}{9.3}\selectfont P04 · Un geyser: associazione e previsione};
\node[anchor=north east,inner sep=0] at ([xshift=17.65cm,yshift=-12.06cm]current page.north west){\color{Classical}\fontsize{7.3}{9.1}\selectfont };
\node[anchor=north east,inner sep=0] at ([xshift=19.05cm,yshift=-12.06cm]current page.north west){\color{Muted}\fontsize{7.7}{9.3}\selectfont 17/21};
\end{tikzpicture}\end{frame}
\begin{lrbox}{\TitleBox}\begin{minipage}[t]{18.2cm}\raggedright\sloppy\color{Ink}\fontsize{20.7}{24}\selectfont \bfseries\color{Navy}Il miglioramento viene quantificato in minuti\end{minipage}\end{lrbox}
\begin{lrbox}{\BodyBox}\begin{minipage}[t]{10.25cm}\raggedright\sloppy\color{Ink}\fontsize{12.4}{16.4}\selectfont Sul campione di verifica, la RMSE della previsione costante è circa 12,768 minuti, mentre quella della retta è circa 5,850 minuti. La MAE passa rispettivamente da circa 11,016 a 4,941 minuti. La durata fornisce quindi informazione predittiva utile entro questo esperimento e la separazione dichiarata. La grandezza del miglioramento rimane riferita a questi eventi e alle condizioni di scambiabilità, mentre la valutazione di utilità operativa richiederebbe soglie di errore rilevanti e dati cronologici nuovi.\end{minipage}\end{lrbox}
\begin{lrbox}{\SideBox}\begin{minipage}[t]{7.45cm}\raggedright\sloppy\color{Ink}\fontsize{12.2}{16.1}\selectfont \includegraphics[width=7.45cm]{figure/P04/F06.pdf}\par\vspace{.16cm}{\fontsize{9.5}{12.6}\selectfont RMSE calcolata sugli stessi 82 eventi. La media di riferimento proviene esclusivamente dal train.}\end{minipage}\end{lrbox}
\begin{lrbox}{\RecallBox}\begin{minipage}[t]{16.45cm}\raggedright\sloppy\color{Ink}\fontsize{10.1}{13.2}\selectfont Il guadagno predittivo viene espresso rispetto a un riferimento e a un dominio di valutazione dichiarati.\end{minipage}\end{lrbox}
\MeasureBox{18-TitleBox}{\TitleBox}
\MeasureBox{18-BodyBox}{\BodyBox}
\MeasureBox{18-SideBox}{\SideBox}
\MeasureBox{18-RecallBox}{\RecallBox}
\begin{frame}[plain]\hypertarget{slide18}{}
\bookmark[level=1,page=18]{Il miglioramento viene quantificato in minuti}
\begin{tikzpicture}[remember picture,overlay]\node[anchor=north east,inner sep=0] at ([xshift=19.05cm,yshift=-0.16cm]current page.north west){\color{Classical}\fontsize{8}{10}\selectfont\hyperlink{slide2}{Indice}};
\node[anchor=north west,inner sep=0] at ([xshift=0.95cm,yshift=-0.62cm]current page.north west){\usebox{\TitleBox}};
\draw[Classical,line width=1.3pt] ([xshift=.95cm,yshift=-1.68cm]current page.north west)--([xshift=2.25cm,yshift=-1.68cm]current page.north west);\node[anchor=north west,inner sep=0] at ([xshift=0.95cm,yshift=-2.12cm]current page.north west){\usebox{\BodyBox}};
\node[anchor=north west,inner sep=0] at ([xshift=11.6cm,yshift=-2.1cm]current page.north west){\color{Quantum}\bfseries\fontsize{11.5}{14}\selectfont Evidenza grafica};
\node[anchor=north west,inner sep=0] at ([xshift=11.6cm,yshift=-2.75cm]current page.north west){\usebox{\SideBox}};
\draw[Rule,line width=.5pt] ([xshift=.95cm,yshift=-10.43cm]current page.north west)--([xshift=19.05cm,yshift=-10.43cm]current page.north west);\node[anchor=north west,inner sep=0] at ([xshift=0.95cm,yshift=-10.66cm]current page.north west){\color{Classical}\bfseries\fontsize{9.4}{12}\selectfont Richiami};
\node[anchor=north west,inner sep=0] at ([xshift=2.6cm,yshift=-10.64cm]current page.north west){\usebox{\RecallBox}};
\draw[Rule,line width=.35pt] ([xshift=.95cm,yshift=-11.92cm]current page.north west)--([xshift=19.05cm,yshift=-11.92cm]current page.north west);\node[anchor=north west,inner sep=0] at ([xshift=0.95cm,yshift=-12.06cm]current page.north west){\color{Muted}\fontsize{7.7}{9.3}\selectfont P04 · Un geyser: associazione e previsione};
\node[anchor=north east,inner sep=0] at ([xshift=17.65cm,yshift=-12.06cm]current page.north west){\color{Classical}\fontsize{7.3}{9.1}\selectfont };
\node[anchor=north east,inner sep=0] at ([xshift=19.05cm,yshift=-12.06cm]current page.north west){\color{Muted}\fontsize{7.7}{9.3}\selectfont 18/21};
\end{tikzpicture}\end{frame}
\begin{lrbox}{\TitleBox}\begin{minipage}[t]{18.2cm}\raggedright\sloppy\color{Ink}\fontsize{20.7}{24}\selectfont \bfseries\color{Navy}La copertura viene osservata nel test\end{minipage}\end{lrbox}
\begin{lrbox}{\BodyBox}\begin{minipage}[t]{10.25cm}\raggedright\sloppy\color{Ink}\fontsize{12.4}{16.4}\selectfont Ottanta delle ottantadue attese ricadono nei rispettivi intervalli predittivi classici al 95\%, producendo una copertura osservata di circa 97,56\%. La figura ordina gli eventi per previsione e rende riconoscibili le due osservazioni esterne alla fascia. Questa proporzione costituisce una verifica empirica su un campione limitato, soggetta a variabilità e alle condizioni della separazione. La vicinanza al livello nominale non garantisce la medesima copertura per eventi futuri acquisiti con protocolli o condizioni del sistema differenti.\par\vspace{.28cm}{\centering\fontsize{13}{17}\selectfont\(\displaystyle\begin{gathered}\widehat C=\frac1{82}\sum_i\mathbf1\{L_i\leq y_i\leq U_i\}=\frac{80}{82}\end{gathered}\)\par}\end{minipage}\end{lrbox}
\begin{lrbox}{\SideBox}\begin{minipage}[t]{7.45cm}\raggedright\sloppy\color{Ink}\fontsize{12.2}{16.1}\selectfont \includegraphics[width=7.45cm]{figure/P04/F08.pdf}\par\vspace{.16cm}{\fontsize{9.5}{12.6}\selectfont $L_i,U_i$: estremi predittivi; $y_i$: attesa osservata. Eventi ordinati per previsione, senza cronologia; ocra: due casi esterni.}\end{minipage}\end{lrbox}
\begin{lrbox}{\RecallBox}\begin{minipage}[t]{16.45cm}\raggedright\sloppy\color{Ink}\fontsize{10.1}{13.2}\selectfont La copertura nominale è una proprietà della procedura sotto ipotesi; quella osservata è una proporzione campionaria.\end{minipage}\end{lrbox}
\MeasureBox{19-TitleBox}{\TitleBox}
\MeasureBox{19-BodyBox}{\BodyBox}
\MeasureBox{19-SideBox}{\SideBox}
\MeasureBox{19-RecallBox}{\RecallBox}
\begin{frame}[plain]\hypertarget{slide19}{}
\bookmark[level=1,page=19]{La copertura viene osservata nel test}
\begin{tikzpicture}[remember picture,overlay]\node[anchor=north east,inner sep=0] at ([xshift=19.05cm,yshift=-0.16cm]current page.north west){\color{Classical}\fontsize{8}{10}\selectfont\hyperlink{slide2}{Indice}};
\node[anchor=north west,inner sep=0] at ([xshift=0.95cm,yshift=-0.62cm]current page.north west){\usebox{\TitleBox}};
\draw[Classical,line width=1.3pt] ([xshift=.95cm,yshift=-1.68cm]current page.north west)--([xshift=2.25cm,yshift=-1.68cm]current page.north west);\node[anchor=north west,inner sep=0] at ([xshift=0.95cm,yshift=-2.12cm]current page.north west){\usebox{\BodyBox}};
\node[anchor=north west,inner sep=0] at ([xshift=11.6cm,yshift=-2.1cm]current page.north west){\color{Quantum}\bfseries\fontsize{11.5}{14}\selectfont Evidenza grafica};
\node[anchor=north west,inner sep=0] at ([xshift=11.6cm,yshift=-2.75cm]current page.north west){\usebox{\SideBox}};
\draw[Rule,line width=.5pt] ([xshift=.95cm,yshift=-10.43cm]current page.north west)--([xshift=19.05cm,yshift=-10.43cm]current page.north west);\node[anchor=north west,inner sep=0] at ([xshift=0.95cm,yshift=-10.66cm]current page.north west){\color{Classical}\bfseries\fontsize{9.4}{12}\selectfont Richiami};
\node[anchor=north west,inner sep=0] at ([xshift=2.6cm,yshift=-10.64cm]current page.north west){\usebox{\RecallBox}};
\draw[Rule,line width=.35pt] ([xshift=.95cm,yshift=-11.92cm]current page.north west)--([xshift=19.05cm,yshift=-11.92cm]current page.north west);\node[anchor=north west,inner sep=0] at ([xshift=0.95cm,yshift=-12.06cm]current page.north west){\color{Muted}\fontsize{7.7}{9.3}\selectfont P04 · Un geyser: associazione e previsione};
\node[anchor=north east,inner sep=0] at ([xshift=17.65cm,yshift=-12.06cm]current page.north west){\color{Classical}\fontsize{7.3}{9.1}\selectfont };
\node[anchor=north east,inner sep=0] at ([xshift=19.05cm,yshift=-12.06cm]current page.north west){\color{Muted}\fontsize{7.7}{9.3}\selectfont 19/21};
\end{tikzpicture}\end{frame}
\begin{lrbox}{\TitleBox}\begin{minipage}[t]{18.2cm}\raggedright\sloppy\color{Ink}\fontsize{20.7}{24}\selectfont \bfseries\color{Navy}Le pendenze interne differiscono da quella globale\end{minipage}\end{lrbox}
\begin{lrbox}{\BodyBox}\begin{minipage}[t]{10.25cm}\raggedright\sloppy\color{Ink}\fontsize{12.4}{16.4}\selectfont Dividendo il solo campione di adattamento alla durata didattica di tre minuti, le pendenze risultano circa 5,588 e 4,328 nei gruppi di 72 e 118 eventi. Entrambe sono inferiori alla pendenza globale, circa 10,898, che incorpora anche la separazione fra addensamenti. Il confronto descrive una sensibilità dell’associazione e non assegna automaticamente uno stato fisico ai due gruppi. La soglia non viene scelta attraverso il test e le rette separate non sostituiscono retroattivamente la procedura già valutata.\end{minipage}\end{lrbox}
\begin{lrbox}{\SideBox}\begin{minipage}[t]{7.45cm}\raggedright\sloppy\color{Ink}\fontsize{12.2}{16.1}\selectfont \includegraphics[width=7.45cm]{figure/P04/F07.pdf}\par\vspace{.16cm}{\fontsize{9.5}{12.6}\selectfont Rette entro due intervalli di durata, sul solo train. La soglia di tre minuti è un criterio descrittivo dichiarato.}\end{minipage}\end{lrbox}
\begin{lrbox}{\RecallBox}\begin{minipage}[t]{16.45cm}\raggedright\sloppy\color{Ink}\fontsize{10.1}{13.2}\selectfont L’associazione complessiva può riflettere sia variazioni interne sia differenze fra gruppi.\end{minipage}\end{lrbox}
\MeasureBox{20-TitleBox}{\TitleBox}
\MeasureBox{20-BodyBox}{\BodyBox}
\MeasureBox{20-SideBox}{\SideBox}
\MeasureBox{20-RecallBox}{\RecallBox}
\begin{frame}[plain]\hypertarget{slide20}{}
\bookmark[level=0,page=20]{Struttura interna e discussione}
\bookmark[level=1,page=20]{Le pendenze interne differiscono da quella globale}
\begin{tikzpicture}[remember picture,overlay]\node[anchor=north east,inner sep=0] at ([xshift=19.05cm,yshift=-0.16cm]current page.north west){\color{Classical}\fontsize{8}{10}\selectfont\hyperlink{slide2}{Indice}};
\node[anchor=north west,inner sep=0] at ([xshift=0.95cm,yshift=-0.62cm]current page.north west){\usebox{\TitleBox}};
\draw[Classical,line width=1.3pt] ([xshift=.95cm,yshift=-1.68cm]current page.north west)--([xshift=2.25cm,yshift=-1.68cm]current page.north west);\node[anchor=north west,inner sep=0] at ([xshift=0.95cm,yshift=-2.12cm]current page.north west){\usebox{\BodyBox}};
\node[anchor=north west,inner sep=0] at ([xshift=11.6cm,yshift=-2.1cm]current page.north west){\color{Quantum}\bfseries\fontsize{11.5}{14}\selectfont Evidenza grafica};
\node[anchor=north west,inner sep=0] at ([xshift=11.6cm,yshift=-2.75cm]current page.north west){\usebox{\SideBox}};
\draw[Rule,line width=.5pt] ([xshift=.95cm,yshift=-10.43cm]current page.north west)--([xshift=19.05cm,yshift=-10.43cm]current page.north west);\node[anchor=north west,inner sep=0] at ([xshift=0.95cm,yshift=-10.66cm]current page.north west){\color{Classical}\bfseries\fontsize{9.4}{12}\selectfont Richiami};
\node[anchor=north west,inner sep=0] at ([xshift=2.6cm,yshift=-10.64cm]current page.north west){\usebox{\RecallBox}};
\draw[Rule,line width=.35pt] ([xshift=.95cm,yshift=-11.92cm]current page.north west)--([xshift=19.05cm,yshift=-11.92cm]current page.north west);\node[anchor=north west,inner sep=0] at ([xshift=0.95cm,yshift=-12.06cm]current page.north west){\color{Muted}\fontsize{7.7}{9.3}\selectfont P04 · Un geyser: associazione e previsione};
\node[anchor=north east,inner sep=0] at ([xshift=17.65cm,yshift=-12.06cm]current page.north west){\color{Classical}\fontsize{7.3}{9.1}\selectfont };
\node[anchor=north east,inner sep=0] at ([xshift=19.05cm,yshift=-12.06cm]current page.north west){\color{Muted}\fontsize{7.7}{9.3}\selectfont 20/21};
\end{tikzpicture}\end{frame}
\begin{lrbox}{\TitleBox}\begin{minipage}[t]{18.2cm}\raggedright\sloppy\color{Ink}\fontsize{20.7}{24}\selectfont \bfseries\color{Navy}Consegna: distinguere effetto medio e previsione\end{minipage}\end{lrbox}
\begin{lrbox}{\BodyBox}\begin{minipage}[t]{10.25cm}\raggedright\sloppy\color{Ink}\fontsize{12.4}{16.4}\selectfont La relazione presenta fonte, separazione, grafici del train, pendenza con intervalli, residui, confronto predittivo e copertura. L’esercizio chiede quale intervallo utilizzare per la media degli eventi a una durata fissata e quale per un singolo nuovo evento, motivando il termine aggiuntivo nella varianza. Una seconda discussione identifica la fuga di informazione introdotta dalla media del test. La conclusione mantiene distinti associazione nel campione, generalizzazione interna e requisiti di una previsione cronologica applicata al sistema reale.\end{minipage}\end{lrbox}
\begin{lrbox}{\SideBox}\begin{minipage}[t]{7.45cm}\raggedright\sloppy\color{Ink}\fontsize{12.2}{16.1}\selectfont R/P04\_figure.R conserva le previsioni per identificativo e i calcoli della covarianza HC3. I dati e il piano di separazione si trovano in dati/progetti. Il progetto raccorda regressione, intervalli predittivi e validazione dei moduli M07, M10 e M16.\end{minipage}\end{lrbox}
\begin{lrbox}{\RecallBox}\begin{minipage}[t]{16.45cm}\raggedright\sloppy\color{Ink}\fontsize{10.1}{13.2}\selectfont La forza della conclusione deriva dal collegamento fra bersaglio, informazione disponibile e disegno della verifica.\end{minipage}\end{lrbox}
\MeasureBox{21-TitleBox}{\TitleBox}
\MeasureBox{21-BodyBox}{\BodyBox}
\MeasureBox{21-SideBox}{\SideBox}
\MeasureBox{21-RecallBox}{\RecallBox}
\begin{frame}[plain]\hypertarget{slide21}{}
\bookmark[level=1,page=21]{Consegna: distinguere effetto medio e previsione}
\begin{tikzpicture}[remember picture,overlay]\node[anchor=north east,inner sep=0] at ([xshift=19.05cm,yshift=-0.16cm]current page.north west){\color{Classical}\fontsize{8}{10}\selectfont\hyperlink{slide2}{Indice}};
\node[anchor=north west,inner sep=0] at ([xshift=0.95cm,yshift=-0.62cm]current page.north west){\usebox{\TitleBox}};
\draw[Classical,line width=1.3pt] ([xshift=.95cm,yshift=-1.68cm]current page.north west)--([xshift=2.25cm,yshift=-1.68cm]current page.north west);\node[anchor=north west,inner sep=0] at ([xshift=0.95cm,yshift=-2.12cm]current page.north west){\usebox{\BodyBox}};
\node[anchor=north west,inner sep=0] at ([xshift=11.6cm,yshift=-2.1cm]current page.north west){\color{Quantum}\bfseries\fontsize{11.5}{14}\selectfont Sviluppo e interpretazione};
\node[anchor=north west,inner sep=0] at ([xshift=11.6cm,yshift=-2.75cm]current page.north west){\usebox{\SideBox}};
\draw[Rule,line width=.5pt] ([xshift=.95cm,yshift=-10.43cm]current page.north west)--([xshift=19.05cm,yshift=-10.43cm]current page.north west);\node[anchor=north west,inner sep=0] at ([xshift=0.95cm,yshift=-10.66cm]current page.north west){\color{Classical}\bfseries\fontsize{9.4}{12}\selectfont Richiami};
\node[anchor=north west,inner sep=0] at ([xshift=2.6cm,yshift=-10.64cm]current page.north west){\usebox{\RecallBox}};
\draw[Rule,line width=.35pt] ([xshift=.95cm,yshift=-11.92cm]current page.north west)--([xshift=19.05cm,yshift=-11.92cm]current page.north west);\node[anchor=north west,inner sep=0] at ([xshift=0.95cm,yshift=-12.06cm]current page.north west){\color{Muted}\fontsize{7.7}{9.3}\selectfont P04 · Un geyser: associazione e previsione};
\node[anchor=north east,inner sep=0] at ([xshift=17.65cm,yshift=-12.06cm]current page.north west){\color{Classical}\fontsize{7.3}{9.1}\selectfont R datasets: faithful; statsmodels, HC3\_se};
\node[anchor=north east,inner sep=0] at ([xshift=19.05cm,yshift=-12.06cm]current page.north west){\color{Muted}\fontsize{7.7}{9.3}\selectfont 21/21};
\end{tikzpicture}\end{frame}
\end{document}